\documentclass[12pt]{article}

\begin{document}

\textbf{Preface}

\hfill

In the 30 years since the publication of the first edition of this book, much
has changed in the field of mathematics known as dynamical systems. In the
early 1970s, we had very little access to high-speed computers and computer
graphics. The word chaos had never been used in a mathematical setting, and
most of the interest in the theory of differential equations and dynamical
systems was confined to a relatively small group of mathematicians.

Things have changed dramatically in the ensuing 3 decades. Computers are
everywhere, and software packages that can be used to approximate solutions
of differential equations and view the results graphically are widely available.
As a consequence, the analysis of nonlinear systems of differential equations
is much more accessible than it once was. The discovery of such complicated
dynamical systems as the horseshoe map, homoclinic tangles, and the
Lorenz system, and their mathematical analyses, convinced scientists that simple
stable motions such as equilibria or periodic solutions were not always the
most important behavior of solutions of differential equations. The beauty
and relative accessibility of these chaotic phenomena motivated scientists and
engineers in many disciplines to look more carefully at the important differential
equations in their own fields. In many cases, they found chaotic behavior in
these systems as well. Now dynamical systems phenomena appear in virtually
every area of science, from the oscillating Belousov-Zhabotinsky reaction in
chemistry to the chaotic Chua circuit in electrical engineering, from complicated
motions in celestial mechanics to the bifurcations arising in ecological
systems.

As a consequence, the audience for a text on differential equations and
dynamical systems is considerably larger and more diverse than it was in
the 1970s. We have accordingly made several major structural changes to
this text, including the following:

\begin{enumerate}
\item The treatment of linear algebra has been scaled back. We have dispensed
with the generalities involved with abstract vector spaces and normed
linear spaces. We no longer include a complete proof of the reduction
of all n × n matrices to canonical form. Rather we deal primarily with
matrices no larger than 4 × 4.
\item We have included a detailed discussion of the chaotic behavior in the
Lorenz attractor, the Shil’nikov system, and the double scroll attractor.
\item Many new applications are included; previous applications have been
updated.
\item There are now several chapters dealing with discrete dynamical systems.
\item We deal primarily with systems that are $C^{\infty}$, thereby simplifying many of
the hypotheses of theorems.
\end{enumerate}

The book consists of three main parts. The first part deals with linear systems
of differential equations together with some first-order nonlinear equations.
The second part of the book is the main part of the text: Here we concentrate on
nonlinear systems, primarily two dimensional, as well as applications of these
systems in a wide variety of fields. The third part deals with higher dimensional
systems. Here we emphasize the types of chaotic behavior that do not occur in
planar systems, as well as the principal means of studying such behavior, the
reduction to a discrete dynamical system.

Writing a book for a diverse audience whose backgrounds vary greatly poses
a significant challenge. We view this book as a text for a second course in
differential equations that is aimed not only at mathematicians, but also at
scientists and engineers who are seeking to develop sufficient mathematical
skills to analyze the types of differential equations that arise in their disciplines.
Many who come to this book will have strong backgrounds in linear algebra
and real analysis, but others will have less exposure to these fields. To make
this text accessible to both groups, we begin with a fairly gentle introduction
to low-dimensional systems of differential equations. Much of this will be a
review for readers with deeper backgrounds in differential equations, so we
intersperse some new topics throughout the early part of the book for these
readers.

For example, the first chapter deals with first-order equations. We begin
this chapter with a discussion of linear differential equations and the logistic
population model, topics that should be familiar to anyone who has a rudimentary
acquaintance with differential equations. Beyond this review, we discuss
the logistic model with harvesting, both constant and periodic. This allows
us to introduce bifurcations at an early stage as well as to describe Poincaré
maps and periodic solutions. These are topics that are not usually found in
elementary differential equations courses, yet they are accessible to anyone
with a background in multivariable calculus. Of course, readers with a limited
background may wish to skip these specialized topics at first and concentrate
on the more elementary material.

Chapters 2 through 6 deal with linear systems of differential equations.
Again we begin slowly, with Chapters 2 and 3 dealing only with planar systems
of differential equations and two-dimensional linear algebra. Chapters
5 and 6 introduce higher dimensional linear systems; however, our emphasis
remains on three- and four-dimensional systems rather than completely
general n-dimensional systems, though many of the techniques we describe
extend easily to higher dimensions.

The core of the book lies in the second part. Here we turn our attention
to nonlinear systems. Unlike linear systems, nonlinear systems present
some serious theoretical difficulties such as existence and uniqueness of solutions,
dependence of solutions on initial conditions and parameters, and the
like. Rather than plunge immediately into these difficult theoretical questions,
which require a solid background in real analysis, we simply state the important
results in Chapter 7 and present a collection of examples that illustrate
what these theorems say (and do not say). Proofs of all of these results are
included in the final chapter of the book.

In the first few chapters in the nonlinear part of the book, we introduce
such important techniques as linearization near equilibria, nullcline analysis,
stability properties, limit sets, and bifurcation theory. In the latter half of this
part, we apply these ideas to a variety of systems that arise in biology, electrical
engineering, mechanics, and other fields.

Many of the chapters conclude with a section called “Exploration.” These
sections consist of a series of questions and numerical investigations dealing
with a particular topic or application relevant to the preceding material. In
each Exploration we give a brief introduction to the topic at hand and provide
references for further reading about this subject. But we leave it to the reader to
tackle the behavior of the resulting system using the material presented earlier.
We often provide a series of introductory problems as well as hints as to how
to proceed, but in many cases, a full analysis of the system could become a
major research project. You will not find “answers in the back of the book” for
these questions; in many cases nobody knows the complete answer. (Except,
of course, you!)

The final part of the book is devoted to the complicated nonlinear behavior
of higher dimensional systems known as chaotic behavior. We introduce these
ideas via the famous Lorenz system of differential equations. As is often the
case in dimensions three and higher, we reduce the problem of comprehending
the complicated behavior of this differential equation to that of understanding
the dynamics of a discrete dynamical system or iterated function. So we then
take a detour into the world of discrete systems, discussing along the way how
symbolic dynamics may be used to describe completely certain chaotic systems.
We then return to nonlinear differential equations to apply these techniques
to other chaotic systems, including those that arise when homoclinic orbits
are present.

We maintain a website at math.bu.edu/hsd devoted to issues regarding
this text. Look here for errata, suggestions, and other topics of interest to
teachers and students of differential equations. We welcome any contributions
from readers at this site.

It is a special pleasure to thank Bard Ermentrout, John Guckenheimer, Tasso
Kaper, Jerrold Marsden, and Gareth Roberts for their many fine comments
about an earlier version of this edition. Thanks are especially due to Daniel
Look and Richard Moeckel for a careful reading of the entire manuscript.
Many of the phase plane drawings in this book were made using the excellent
Mathematica package called DynPac: A Dynamical Systems Package for Mathematica
written by Al Clark. See www.me.rochester.edu/~clark/dynpac.html.
And, as always, Killer Devaney digested the entire manuscript; all errors that
remain are due to her.

\begin{thebibliography}{9}
\bibitem{texbook}
Morris W. Hirsch, Stephen Smale, Robert L. Devaney (2004) \emph{DIFFERENTIAL EQUATIONS,
DYNAMICAL SYSTEMS, AND
AN INTRODUCTION
TO CHAOS}, ELSEVIER.

\bibitem{lamport94}
Leslie Lamport (1994) \emph{\LaTeX: a document preparation system}, Addison
Wesley, Massachusetts, 2nd ed.
\end{thebibliography}


\end{document}
